\iflanguage{ngerman}
{\chapter{Verwandte Arbeiten}}
{\chapter{Related Work}}
\label{sec:related}

\section{Introduction to Clinical Predictive Modeling and EHR Data}
\subsection{The Critical Role of Automated Diagnosis Prediction in Modern Healthcare}
Automated diagnosis prediction represents a fundamental advancement in health informatics, offering substantial benefits across the continuum of patient care. The ability to accurately predict diagnoses, whether retrospectively for clinical coding or proactively for patient forecasting, is essential for refining clinical decision-making, optimizing institutional resource allocation, and ultimately improving patient outcomes.1The scope of predictive modeling spans two distinct, yet interconnected, objectives: determining the diagnoses concurrent with the current hospital visit and forecasting diagnoses that may arise in a future temporal window. These tasks address different clinical requirements, ranging from administrative efficiency to strategic preventative care planning.
The integration of artificial intelligence (AI) into healthcare workflows is increasingly dependent on the quality and comprehensiveness of the data used. Data-driven deep learning methodologies have achieved considerable advancements in healthcare practices, particularly in predicting next admissions and subsequent diagnoses.1This systematic leverage of accumulated Electronic Health Record (EHR) data is transforming the capacity of healthcare providers to anticipate patient needs and implement proactive future plans.

\subsection{The Multimodal and Sequential Challenges of Electronic Health Records (EHRs)}
EHR data presents significant modeling challenges due to its inherent complexity, characterized by both heterogeneity and sequential dependency. EHRs encompass a range of modalities, broadly categorized into structured data—including medical codes (ICD, LOINC, ATC), demographics, and time series—and unstructured data, primarily consisting of free-text clinical notes, such as discharge summaries and progress notes.1Extracting meaningful insights requires models capable of integrating these disparate sources effectively.4
A persistent limitation noted in many preceding predictive studies is the failure to adequately address the heterogeneous and hierarchical characteristics intrinsic to EHR data.1The lack of sophisticated modeling for these features inevitably leads to suboptimal predictive performance. Therefore, effective models must not only fuse multimodal information but also capture the semantic relationships between medical codes (hierarchy) and the chronological sequence of clinical events (temporal dynamics).5Furthermore, the MIMIC dataset, providing extensive, de-identified longitudinal patient information from intensive care units, serves as a crucial and widely adopted foundation for developing and validating advanced language models and predictive analytics in this domain.6
A key requirement for any longitudinal prediction task is the rigorous enforcement of temporal causality. The ability of a model to accurately predict a future event, such as a "Next Visit Diagnosis," is entirely dependent upon its capacity to reason about sequential clinical events and dependencies. Conventional encoding approaches often fail to capture these rich temporal dynamics, which necessitates specialized architectural designs.5If the sequence model architecture neglects to enforce a causal (unidirectional) masking mechanism or if the data serialization ignores strict timestamps, the integrity of the longitudinal prediction task is compromised by the implicit risk of data leakage (or look-ahead bias) from future events.9The success of models fine-tuned on autoregressive tasks, such as the Next Event Prediction (NEP) paradigm, validates the focus on explicitly modeling disease progression patterns, which is non-trivial for complex, sparse EHR sequences.


\section{Evolution of Deep Learning for EHR Sequence Modeling}
\subsection{Historical Deep Learning Methods: Recurrent Architectures}
The early integration of deep learning (DL) into healthcare predictive modeling was largely dominated by Recurrent Neural Networks (RNNs) and their derivatives, such as Long Short-Term Memory (LSTM) and Gated Recurrent Units (GRU). These architectures proved valuable for processing sequences of diagnostic codes, focusing on representing patient medical history through sequences of visit embeddings.10Their recursive nature allowed them to model temporal dependencies, recognizing that a current diagnosis is often contingent upon a patient's preceding medical history. This approach was instrumental in demonstrating the predictive value of sequential diagnostic information for various patient outcomes.

\subsection{The Transition to Transformer Architectures}
While RNNs established the field, they are inherently limited in capturing very long-range dependencies, a necessity in longitudinal EHR modeling where the time elapsed between visits can span months or years. Consequently, the field has seen a significant shift toward Transformer architectures, which now account for approximately 26\% of DL-based models in this area.10Transformers leverage sophisticated attention mechanisms to weigh the importance of all preceding clinical events, effectively overcoming the temporal bottleneck of RNNs.
However, earlier Transformer models designed for this domain, such as BEHRT, often focused solely on forecasting disorders occurring in the next patient hospital visit or within a specific predefined timeframe (eg, 6 or 12 months).3Furthermore, these models typically employed a multilabel approach, which becomes computationally and conceptually challenging as the universe of forecasted biomedical concepts—such as diagnosis codes—expands.3The complexity of capturing long-term sequences and the sheer scale of the output space necessitated architectures capable of both fine-grained temporal reasoning and large-scale output generation.

\subsection{Multimodal Integration and Feature Fusion Strategies}
Contemporary predictive modeling recognizes that structured diagnostic sequences alone are insufficient. Patient status is holistically represented by combining structured codes with clinical narratives. Empirical evidence consistently confirms that models achieving data fusion by combining structured and unstructured data learn superior patient representations, thereby improving prediction performance and reducing errors.4
The multimodal integration strategy often involves presenting the input features as sequences of visit embeddings, where diagnosis codes are augmented by other clinical features.10Medications, typically represented by codes like ATC or RxNorm, are recognized as the most common additional structured feature integrated alongside diagnosis sequences, appearing in 45\% of relevant studies.10This practice underscores the importance of the selected multimodal input set (medications (ATC), lab tests (LOINC), and procedural codes (ICD-CM)) for comprehensive patient state representation.


\section{Analysis of Bimodal Prediction Tasks: Concurrent vs. Longitudinal}
The predictive tasks investigated fall into two distinct categories: Current Visit Diagnosis Prediction and Next Visit Diagnosis Prediction. While both utilize the same foundational EHR data, they impose vastly different requirements on the model's architecture and contextual reasoning capabilities.

\subsection{Current Visit Diagnosis Prediction: : Concurrent Context Synthesis}
The Current Visit Diagnosis task focuses on generating the final ICD-9-CM diagnosis codes for a single hospital admission. The input context comprises all available data from that visit, including structured codes (medications, labs, procedures) and crucial parsed sections from the discharge summary, such as "Chief complaint," "Brief hospital course," "History of Present Illness (HPI)," and "Family History".11
This task is fundamentally one of concurrent context synthesis. Effective models must integrate information streams available during or at the conclusion of the admission. This necessitates architectural strategies that employ cross-modal attention and sophisticated feature pooling to dynamically adjust to the shifting significance of textual versus tabular evidence.12The performance is often critical during the crucial early stages of a hospital stay when evidence is limited.12The reliance on textual information, specifically detailed clinical notes, is a strong determinant of accuracy in predicting the corresponding diagnosis codes.

\subsection{Next Visit Diagnosis Prediction: Causal Temporal Forecasting}
The Next Visit Diagnosis task represents a true temporal forecasting challenge, requiring the prediction of diagnoses that occur within a distant future window (typically 6 or 12 months).1The input must include the current visit diagnosis alongside the complete preceding longitudinal history.
The cutting edge of sequential modeling frames this challenge using the Next Event Prediction (NEP) paradigm.5Within this framework, EHR data is meticulously serialized into timestamped event chains, capturing the inherent sequential and temporal nature of patient progression, unlike models that rely on static or aggregated data points.9The core methodology involves fine-tuning autoregressive, decoder-based Large Language Models (LLMs)—such as the Llama series—to predict the future clinical event given the historical sequence.5Crucially, the training paradigm must employ a causal (unidirectional) attention mask, which prevents data leakage and ensures the model explicitly models the flow of disease progression, leading to robust representations validated by attention weights that align with clinical pathways.


\section{Contrasting Requirements and Architectural Alignment}
The two tasks place a differential burden on the LLM's parameter capacity, especially when integrating multiple knowledge sources. The success of smaller models, such as the 1-billion parameter Llama variant, in the 'Next Visit' task will serve as a critical test of generalizability and robustness. Temporal forecasting requires deep relational memory and the ability to synthesize sparse, chronologically distant events. If external augmentation methods, such as Knowledge Graphs (KGs) and Retrieval-Augmented Generation (RAG), fail to yield significant gains for the smallest model compared to its 8-billion parameter counterpart, it would suggest that the foundational capacity to process and utilize high-dimensional context is the primary bottleneck, irrespective of the external knowledge provided. This scenario emphasizes that there is a lower boundary on model size required for handling the inherent complexity of longitudinal, multimodal EHR data.
The contrasting requirements of these tasks are summarized below:
\begin{table}[htbp]
\centering
\caption{Comparison of Diagnosis Prediction Tasks in EHR Literature}
\label{tab:task_comparison}
\rowcolors{2}{gray!10}{white}
\begin{tabularx}{0.98\textwidth}{>{\raggedright\arraybackslash}p{2.7cm} 
                                  >{\raggedright\arraybackslash}X 
                                  >{\raggedright\arraybackslash}X}
\toprule
\textbf{Feature} & \textbf{Current Visit Diagnosis Prediction} & \textbf{Next Visit Diagnosis Prediction (Temporal Forecasting)} \\
\midrule
Primary Goal & Retrospective or concurrent clinical coding/classification. & Forecasting future medical events, typically 6 or 12 months out.\cite{ref3} \\
Input Dependency & Heavily reliant on most recent visit context (structured codes and detailed notes).\cite{ref11} & Heavily reliant on the sequence of prior visits (longitudinal history).\cite{ref1} \\
Key Architectural Focus & Multimodal feature fusion and cross-modal attention for real-time evidence synthesis.\cite{ref4} & Temporal reasoning, causal masking, and explicit sequence representation (e.g., NEP).\cite{ref5} \\
Challenges & High label imbalance and large output space for ICD codes; synthesizing immediate multimodal evidence. & Capturing long-term temporal dependencies and predicting rare future events; modeling disease progression.\cite{ref3} \\
\bottomrule
\end{tabularx}
\end{table}

\section{Large Language Models (LLMs) and Efficient Domain Adaptation}
\subsection{LLM Selection and Generative Capabilities in Medicine}
The adoption of Large Language Models (LLMs) marks a transformative shift, offering generative solutions for complex medical NLP tasks. Decoder-only architectures, exemplified by the Llama series (Llama 3, Llama 3.1), are inherently suited for the generative tasks of code or free-text diagnosis output.13These models have demonstrated remarkable capacity for understanding and generating complex medical text.14Systematic evaluations confirm the high efficacy of these models in extracting and processing structured information from unstructured medical reports.13However, performance variability exists across models and data types, highlighting the necessity of rigorous domain-specific fine-tuning for specialized clinical applications.13
The foundation for this fine-tuning is typically large, publicly available medical datasets. The MIMIC series (MIMIC-III and MIMIC-IV), with records from tens of thousands of intensive care patients, provides the essential corpus of clinical notes and coded data required for specialized LLM adaptation.

\subsection{Low-resource fine-tuning techniques}
Effective fine-tuning of large pre-trained models on specialized clinical datasets mandates the use of computationally efficient methods. Techniques such as Low-Rank Adaptation (LoRA) and its quantized variant (QLoRA, often using 4-bit precision) are standard for efficiently adapting the LLM's existing strengths to specific tasks while ensuring time and memory efficient utilization.14These Parameter-Efficient Fine-Tuning (PEFT) methods minimize the computational and memory footprint, making it feasible to adapt models like the Llama-3.1-8B series.15
Furthermore, adapting the LLM to the specialized language of clinical coding often involves instruction tuning. This process leverages the conversion of target ICD codes into their corresponding descriptive texts. These natural language descriptions are then paired with the input medical notes (eg, discharge summaries) to form a training dataset.15Instruction tuning is a validated method for significantly enhancing a language model's capabilities for specialized domain tasks, ensuring that the model understands the required output format and clinical context.

\subsection{The Formulation of EHR Input for LLM Processing}
To leverage the power of decoder-only LLMs for sequential prediction, the disparate EHR modalities must be correctly serialized into a single chronological token sequence. This sequence must integrate the extracted unstructured text segments (eg, "Chief complaint," "HPI") alongside the structured codes (ATC, LOINC, ICD-Proc).9
In the context of autoregressive prediction, the inherent event frequency imbalances in real-world EHRs present a challenge. For instance, routine lab tests occur far more frequently than rare diagnostic events. Standard training on such a skewed distribution can bias the model toward common events. Advanced methodologies, such as temperature-controlled multinomial sampling, are employed during serialization to counter this imbalance. This technique ensures a more balanced representation of event types throughout the training sequence while preserving the critical chronological order, thus improving the model's ability to predict rarer, yet clinically significant, diagnoses in the 'Next Visit' task.

\section{Knowledge Graph Augmentation and Path Mining}
\subsection{Constructing the Clinical Knowledge Foundation (UMLS)}
External knowledge augmentation via Knowledge Graphs (KGs) represents a crucial strategy for enhancing the reliability and reasoning capacity of LLMs in the clinical domain. The construction of the KG foundation relies on integrating disparate medical terminologies, including ICD9CM, LOINC, ATC, and SNOMED CT, often facilitated through the Unified Medical Language System (UMLS).17UMLS CUIs (Concept Unique Identifiers) act as the universal semantic backbone, enabling the establishment of relationships between clinical codes across different vocabularies, even where direct mapping is not strictly 1:1, such as between UMLS CUIs and ATC-5 codes.18
By leveraging this structured ontological information, KGs mitigate the risks inherent in LLMs, which might otherwise struggle with complex, contextually relevant medical inference. Integrating structured knowledge from KGs enhances diagnostic reasoning, minimizing the risk of diagnostic errors and supporting the quest for trustworthy AI in healthcare.

\subsection{Methodology for Knowledge Path Mining (H1 and H2)}
The KG is constructed to capture the geometric structure and explicit relationships existing between clinical concepts (diagnoses, treatments, and tests).7Once the KG is built from UMLS and other terminologies, specific knowledge paths relevant to a patient's current clinical context are mined.
The strategy involves mining paths of varying length:
\begin{itemize}
    \item H1 Paths (Single-Hop): Capture direct semantic relationships between two clinical concepts (eg, a specific drug treats a specific disease).
    \item H2 Paths (Two-Hop): Capture mediated relationships, providing indirect but potentially crucial contextual links (eg, Drug X treats Disease Y, and Disease Y is associated with Symptom Z).
\end{itemize}

These mined paths are then serialized and provided as contextual prompts to the fine-tuned LLMs.19However, simply providing all possible paths is often counterproductive. Research emphasizes that the true performance benefit of KG integration arises from retrieving and ranking only those knowledge paths that are clinically pertinent and contextually aligned with the patient's specific information.19Techniques often involve graph neural networks (GNNs) and attention-based rankers to identify the most relevant topological information from the KG.19
The deployment of KG augmentation introduces a delicate trade-off regarding the LLM's finite context window. Knowledge paths, particularly the longer H2 paths, promptly consume token space. If the retrieved paths are high in noise or low in clinical relevance, they risk displacing crucial, raw EHR input tokens. The success of KG augmentation, especially the combination of H1 and H2 paths, is therefore contingent upon the retrieval system's precision. If the LLM's performance does not increase significantly with the addition of paths, it suggests that the external context is either redundant or acts as distracting noise, failing to overcome the token cost penalty associated with context window constraints.

\subsection{Ablation Study Design: Decomposing Knowledge Contribution}
The planned ablation study—comparing the baseline model against models augmented with H1 paths, H2 paths, and both H1 and H2 paths—is a rigorous methodological approach designed to decompose the contribution of knowledge complexity.
The comparison between the baseline and the H1-only condition isolates the effect of providing immediate, explicit semantic neighbors.
The relative metric change observed between the H1-only condition and the H2-only condition offers direct evidence regarding the LLM's intrinsic capacity to utilize complex, mediated, or indirect associations. A significant gain from H2 augmentation implies that the model possesses the requisite capability for deep, two-step diagnostic reasoning over abstract knowledge structures.

\section*{Advanced Contextual Enhancement: Retrieval-Augmented Generation (RAG)}
\subsection{The Role of RAG in Clinical Generative Models}
Retrieval-Augmented Generation (RAG) is a powerful technique for grounding the outputs of LLMs, directly addressing challenges such as factually incorrect outputs (hallucinations) and reliance on outdated internal knowledge.13By dynamically enriching the input context with information retrieved from an external, verified database, RAG significantly enhances data accessibility and supports more accurate clinical decision-making processes.13In the safety-critical domain of diagnosis prediction, RAG is instrumental in moving LLM predictions toward greater trustworthiness and reliability.

\subsection{Hybrid RAG Architecture for Clinical History}
The proposed RAG system employs a hybrid context formulation that combines two distinct forms of clinical knowledge:

$$\text{Context} = (\text{Mined H1/H2 Knowledge Paths}) + (\text{Past Cases from Training Sample})$$
\begin{itemize}
    \item \textbf{Abstract, Ontological Knowledge (KG Paths)}: The H1 and H2 paths retrieved from UMLS provide generalized, codified rules and relationships derived from established medical ontologies (eg, X is-a Y, Drug A treats Condition B).
    \item \textbf{Empirical, Case-Specific Knowledge (Past Cases)}: The retrieval of similar past patient cases from the training data provides concrete, real-world examples of disease progression, co-morbidities, and temporal sequences observed within the specific patient population (MIMIC). This empirical knowledge is essential for prediction because it captures local health patterns and observed clinical trajectories that may not be fully formalized in generalized KGs
\end{itemize}
This fusion mechanism ensures that the model benefits from both the abstract rules of medical science and the observed patterns of local clinical practice.

\subsection{The RAG Ablation: Quantifying Empirical Value}
The final ablation step, which compares the performance of the KG-augmented model ($\text{H1} + \text{H2}$) against the hybrid RAG-augmented model ($\text{H1} + \text{H2} + \text{RAG}$), is designed to rigorously quantify the incremental value of empirical patient history.
If the addition of past, similar patient records (RAG component) results in a substantial gain in prediction metrics, particularly for the longitudinal "Next Visit" task, it provides strong evidence. This gain would suggest that accurate disease forecasting relies on specific, local patterns of sequential events and co-occurrence that are derived from population data, rather than being solely deducible from codified relationships in general terminologies. The explicit inclusion of these empirical cases allows the LLM to ground its generative forecast in verifiable, patient-specific precedents.


% \section{Clinical Decision Support and Learning from EHR Data}
% Clinical decision support systems (CDSS) aim to provide patient-specific knowledge and recommendations to clinicians at the point of care, improving safety, quality, and efficiency. Modern CDSS increasingly leverages machine learning models trained on electronic health records (EHRs), where the central challenge is converting heterogeneous clinical data into representations that support robust prediction under missingness, noise, and strong label imbalance \cite{sutton2020cdss,rajkomar2018deepehr}.

% EHRs are inherently multimodal. They combine structured elements (e.g., diagnosis/procedure codes, medication codes, laboratory results, demographics) with unstructured narrative text (e.g., discharge summaries, progress notes). The MIMIC family of datasets is widely used for developing and benchmarking these models, as it provides large-scale, de-identified ICU data including both structured tables and clinical notes \cite{johnson2016mimiciii,johnson2023mimiciv}. The two tasks in this thesis---(\textit{i}) automated ICD coding for an admission and (\textit{ii}) next-visit diagnosis prediction---map naturally to two broad paradigms in the literature: (1) \emph{visit-level multi-label coding} and (2) \emph{longitudinal sequence forecasting}.

% \section{Automated ICD Coding from Clinical Text}
% Automated ICD coding is typically formulated as a multi-label prediction problem: given a clinical document (often the discharge summary), predict a set of ICD codes assigned at discharge. Early work explored classical feature-based approaches and emphasized the importance of evaluation protocols under extreme label imbalance \cite{perotte2014diagnosis}. As datasets grew and neural NLP matured, deep learning models became dominant, largely due to their ability to learn task-specific representations from raw text while scaling to thousands of labels.

% \subsection{CNN and Attention-Based Neural Baselines}
% A major milestone is the Convolutional Attention-based Multi-Label model (CAML), which uses convolutional encoders and label-wise attention to identify text spans most relevant to each code; it also popularized commonly reported metrics such as micro-F1 and precision@k for large-label ICD coding \cite{mullenbach2018caml}. Subsequent convolutional variants improved long-document modeling by using multiple filter widths and residual connections, such as MultiResCNN \cite{li2020multirescnn}. These architectures remain strong baselines because they are computationally efficient and handle long notes more directly than many vanilla Transformer pipelines.

% \subsection{Label-Wise Attention and Hierarchy-Aware Modeling}
% Beyond document encoders, a key design axis in ICD coding is \emph{label-aware} representation learning. LAAT (Label Attention) learns a label-specific document representation and additionally explores exploiting ICD hierarchy to improve prediction for infrequent codes \cite{vu2020laat}. Hierarchical label-wise attention mechanisms have also been used with Transformer backbones to improve interpretability and performance; HiLAT is one such approach, evaluated on common benchmark setups including the top-50 most frequent ICD codes \cite{liu2022hilat}.

% \subsection{Pretrained Language Models for ICD Coding}
% While pretrained language models (PLMs) dominate many NLP tasks, straightforward fine-tuning can underperform in ICD coding due to three recurring issues: very long clinical documents, very large output spaces, and domain mismatch \cite{huang2022plmicd}. PLM-ICD analyzes these failure modes and proposes a more effective PLM-based framework, reporting competitive performance on MIMIC benchmarks \cite{huang2022plmicd}. This line of work motivates careful input structuring, length handling, and training strategies when using Transformer-based or LLM-based architectures for clinical coding.

% \section{Multimodal ICD Coding: Structured EHR + Clinical Text}
% A growing body of work shows that adding structured EHR features can improve ICD coding beyond text-only models, especially when structured signals capture decisive evidence not reliably expressed in narrative notes. Xu et al.\ propose a multimodal framework that trains modality-specific models and combines them via ensembling to predict ICD codes \cite{xu2019multimodal}. More recent methods learn \emph{fusion} within a unified architecture: TreeMAN constructs tree-based features from structured data and uses attention to enhance text representations for multimodal ICD prediction \cite{liu2022treeman}. These approaches are closely aligned with your thesis setting (structured codes + clinical text $\rightarrow$ ICD coding), and they provide precedent for expecting gains from multimodal fusion when the structured channel is informative and complementary.

% \section{Longitudinal Modeling and Next-Visit Diagnosis Prediction}
% Next-visit diagnosis prediction (and related forecasting tasks) requires explicitly modeling temporal progression: given past visits (codes, labs, treatments, notes), predict future diagnoses for the next encounter or a future window. Historically, recurrent neural networks were widely used for such forecasting on sequences of coded events. Doctor AI models longitudinal records with RNNs to predict subsequent diagnoses/medications, demonstrating early success on large-scale visit sequences \cite{choi2016doctorai}. RETAIN introduces reverse-time attention to improve interpretability while maintaining predictive performance \cite{choi2016retain}. Representation learning approaches such as Deep Patient also show that learned patient embeddings can improve downstream predictive tasks \cite{miotto2016deeppatient}.

% \subsection{Transformers for Structured EHR Sequences}
% Transformer models have also been adapted to structured EHR sequences to improve long-range dependency modeling. BEHRT applies a Transformer to EHR event sequences and predicts diagnoses in future visits \cite{li2020behrt}. Med-BERT demonstrates that BERT-style pretraining can be applied to large-scale structured EHR data and can improve downstream disease prediction, especially in lower-resource fine-tuning regimes \cite{rasmy2021medbert}. These works support the broader thesis motivation that \emph{sequence-aware} modeling is crucial for next-visit prediction, and that pretraining/adaptation can help when learning from limited labeled cohorts.

% \section{Large Language Models and Efficient Adaptation for Clinical NLP}
% Modern LLM pipelines build on the Transformer architecture \cite{vaswani2017transformer}. General-purpose LLMs (e.g., LLaMA-family models) provide strong generative capabilities, but adapting them to clinical domains and long-context EHR inputs can be computationally expensive \cite{touvron2023llama}. Parameter-efficient fine-tuning (PEFT) is therefore common: LoRA injects trainable low-rank adapters while freezing base weights \cite{hu2021lora}, and QLoRA combines quantization with LoRA-style adapters to enable memory-efficient fine-tuning of larger models \cite{dettmers2023qlora}. These techniques are especially relevant in clinical modeling settings where GPU resources and training budgets are constrained.

% \section{Knowledge Graphs and Retrieval-Augmented Generation with Medical Facts}
% \subsection{Clinical Knowledge Graphs and UMLS Grounding}
% A central limitation in purely data-driven models is that they may fail on rare codes or clinically complex cases where external biomedical knowledge is required. Medical knowledge graphs (KGs) address this by encoding relationships among clinical concepts (e.g., disease--symptom, drug--indication). Many KG constructions rely on UMLS as a unifying substrate that maps concepts across terminologies via Concept Unique Identifiers (CUIs) \cite{bodenreider2004umls}. In ICD-centric tasks, KGs can provide relational signals that are not explicitly present in an individual patient record and can support better generalization to tail concepts.

% \subsection{RAG as Prompt-Time Knowledge Injection (KG Fact Retrieval)}
% Retrieval-Augmented Generation (RAG) was introduced as a general framework to ground generation in retrieved evidence from an external index, rather than relying only on parametric memory \cite{lewis2020rag}. In healthcare, RAG-style systems are actively explored as a way to improve factuality and provide provenance in decision support settings \cite{ferber2024nejmai,ke2025rag}. 

% This ``KG-facts-in-the-prompt'' pattern has been studied in KG+LLM integrations for clinical reasoning and diagnosis-oriented tasks, where retrieval selects a small, context-relevant subgraph (or paths) to augment the model input \cite{gao2025drknows}. Related work in medical knowledge retrieval also explores strategies for selecting and formatting medical knowledge items for generation, emphasizing the trade-off between evidence coverage and context-window cost \cite{shi2025mkrag}. These findings motivate careful retrieval/reranking (to reduce noise) and controlled prompting (to preserve token budget for patient-specific evidence), both of which are directly relevant to your KG enrichment ablations.

% \section{Evaluation Protocols and Benchmark Variants}
% \subsection{Metrics under Extreme Label Imbalance}
% ICD coding and diagnosis forecasting are highly imbalanced multi-label problems: a small subset of codes occurs frequently, while the majority are rare. Micro-F1 aggregates contributions across all label instances and is dominated by frequent codes; macro-F1 averages F1 per label and is more sensitive to tail performance \cite{perotte2014diagnosis,mullenbach2018caml}. Many ICD coding papers therefore report micro-F1 and macro-F1 together, sometimes complemented by precision@k (e.g., precision@8) to reflect practical coding workflows \cite{mullenbach2018caml}.

% \subsection{Top-$k$ Label Subsets vs. Full-Code Evaluation}
% A common benchmarking variant is to restrict evaluation to the top-$k$ most frequent ICD codes (e.g., top-50) to simplify the label space and stabilize comparisons across methods \cite{mullenbach2018caml,liu2022hilat}. While this setting is useful for ablations and controlled analysis, it is not directly comparable to full-code evaluations that include thousands of labels, where rare-code behavior and calibration become more critical \cite{huang2022plmicd}. Accordingly, studies may not be directly comparable unless they align on (i) the exact label set, (ii) the note types and preprocessing, and (iii) the metric definition (e.g., label set used for macro averaging). This motivates explicit reporting of the label set used to compute micro/macro scores and, when possible, complementary analyses on both frequent and rare subsets.
